% _preamble.tex --- reglages communs IEEE pour tous les TP du parcours "Improve"
% A inclure via % _preamble.tex --- reglages communs IEEE pour tous les TP du parcours "Improve"
% A inclure via % _preamble.tex --- reglages communs IEEE pour tous les TP du parcours "Improve"
% A inclure via % _preamble.tex --- reglages communs IEEE pour tous les TP du parcours "Improve"
% A inclure via \input{_preamble} JUSTE APRES \documentclass[10pt,journal]{IEEEtran}
% Compilation : tectonic (telecharge les packages automatiquement).
\usepackage[T1]{fontenc}
\usepackage[utf8]{inputenc}
\usepackage[french]{babel}
\usepackage{amsmath,amssymb}
\usepackage{newtxtext,newtxmath}     % rendu Times (proche de Tinos) ; APRES amssymb
\usepackage{graphicx}
\usepackage{booktabs}
\usepackage{array}
\usepackage{xcolor}
\usepackage{listings}
\usepackage{enumitem}
\usepackage{tikz}
\usepackage{pgfplots}
\pgfplotsset{compat=1.18}
\usetikzlibrary{arrows.meta,patterns,calc,decorations.pathmorphing}
\usepackage{cite}
\usepackage[hidelinks]{hyperref}

% --- Noms francais des flottants ---
\renewcommand{\tablename}{Tableau}
\renewcommand{\figurename}{Fig.}
\renewcommand{\lstlistingname}{Listing}
\addto\captionsfrench{\renewcommand{\refname}{Références}}

% --- En-tete de revue commun a tous les TP ---
\newcommand{\improvejournal}{IEEE Transactions on Computational Mechanics and Machine Learning, Vol.~1, No.~1, Juin~2026}

% --- Style code Python (mono, sobre facon IEEE) ---
\definecolor{kw}{HTML}{0033B3}
\definecolor{cm}{HTML}{7A7A7A}
\definecolor{st}{HTML}{067D17}
\lstdefinestyle{py}{%
  language=Python, basicstyle=\ttfamily\footnotesize,
  keywordstyle=\color{kw}, commentstyle=\color{cm}\itshape, stringstyle=\color{st},
  numbers=none, showstringspaces=false, breaklines=true, columns=fullflexible,
  aboveskip=3pt, belowskip=1pt, xleftmargin=2pt,
  literate={é}{{\'e}}1 {è}{{\`e}}1 {à}{{\`a}}1 {ê}{{\^e}}1 {ô}{{\^o}}1 {î}{{\^i}}1 {ç}{{\c c}}1 {É}{{\'E}}1 {_}{{\_}}1}
\lstset{style=py}
 JUSTE APRES \documentclass[10pt,journal]{IEEEtran}
% Compilation : tectonic (telecharge les packages automatiquement).
\usepackage[T1]{fontenc}
\usepackage[utf8]{inputenc}
\usepackage[french]{babel}
\usepackage{amsmath,amssymb}
\usepackage{newtxtext,newtxmath}     % rendu Times (proche de Tinos) ; APRES amssymb
\usepackage{graphicx}
\usepackage{booktabs}
\usepackage{array}
\usepackage{xcolor}
\usepackage{listings}
\usepackage{enumitem}
\usepackage{tikz}
\usepackage{pgfplots}
\pgfplotsset{compat=1.18}
\usetikzlibrary{arrows.meta,patterns,calc,decorations.pathmorphing}
\usepackage{cite}
\usepackage[hidelinks]{hyperref}

% --- Noms francais des flottants ---
\renewcommand{\tablename}{Tableau}
\renewcommand{\figurename}{Fig.}
\renewcommand{\lstlistingname}{Listing}
\addto\captionsfrench{\renewcommand{\refname}{Références}}

% --- En-tete de revue commun a tous les TP ---
\newcommand{\improvejournal}{IEEE Transactions on Computational Mechanics and Machine Learning, Vol.~1, No.~1, Juin~2026}

% --- Style code Python (mono, sobre facon IEEE) ---
\definecolor{kw}{HTML}{0033B3}
\definecolor{cm}{HTML}{7A7A7A}
\definecolor{st}{HTML}{067D17}
\lstdefinestyle{py}{%
  language=Python, basicstyle=\ttfamily\footnotesize,
  keywordstyle=\color{kw}, commentstyle=\color{cm}\itshape, stringstyle=\color{st},
  numbers=none, showstringspaces=false, breaklines=true, columns=fullflexible,
  aboveskip=3pt, belowskip=1pt, xleftmargin=2pt,
  literate={é}{{\'e}}1 {è}{{\`e}}1 {à}{{\`a}}1 {ê}{{\^e}}1 {ô}{{\^o}}1 {î}{{\^i}}1 {ç}{{\c c}}1 {É}{{\'E}}1 {_}{{\_}}1}
\lstset{style=py}
 JUSTE APRES \documentclass[10pt,journal]{IEEEtran}
% Compilation : tectonic (telecharge les packages automatiquement).
\usepackage[T1]{fontenc}
\usepackage[utf8]{inputenc}
\usepackage[french]{babel}
\usepackage{amsmath,amssymb}
\usepackage{newtxtext,newtxmath}     % rendu Times (proche de Tinos) ; APRES amssymb
\usepackage{graphicx}
\usepackage{booktabs}
\usepackage{array}
\usepackage{xcolor}
\usepackage{listings}
\usepackage{enumitem}
\usepackage{tikz}
\usepackage{pgfplots}
\pgfplotsset{compat=1.18}
\usetikzlibrary{arrows.meta,patterns,calc,decorations.pathmorphing}
\usepackage{cite}
\usepackage[hidelinks]{hyperref}

% --- Noms francais des flottants ---
\renewcommand{\tablename}{Tableau}
\renewcommand{\figurename}{Fig.}
\renewcommand{\lstlistingname}{Listing}
\addto\captionsfrench{\renewcommand{\refname}{Références}}

% --- En-tete de revue commun a tous les TP ---
\newcommand{\improvejournal}{IEEE Transactions on Computational Mechanics and Machine Learning, Vol.~1, No.~1, Juin~2026}

% --- Style code Python (mono, sobre facon IEEE) ---
\definecolor{kw}{HTML}{0033B3}
\definecolor{cm}{HTML}{7A7A7A}
\definecolor{st}{HTML}{067D17}
\lstdefinestyle{py}{%
  language=Python, basicstyle=\ttfamily\footnotesize,
  keywordstyle=\color{kw}, commentstyle=\color{cm}\itshape, stringstyle=\color{st},
  numbers=none, showstringspaces=false, breaklines=true, columns=fullflexible,
  aboveskip=3pt, belowskip=1pt, xleftmargin=2pt,
  literate={é}{{\'e}}1 {è}{{\`e}}1 {à}{{\`a}}1 {ê}{{\^e}}1 {ô}{{\^o}}1 {î}{{\^i}}1 {ç}{{\c c}}1 {É}{{\'E}}1 {_}{{\_}}1}
\lstset{style=py}
 JUSTE APRES \documentclass[10pt,journal]{IEEEtran}
% Compilation : tectonic (telecharge les packages automatiquement).
\usepackage[T1]{fontenc}
\usepackage[utf8]{inputenc}
\usepackage[french]{babel}
\usepackage{amsmath,amssymb}
\usepackage{newtxtext,newtxmath}     % rendu Times (proche de Tinos) ; APRES amssymb
\usepackage{graphicx}
\usepackage{booktabs}
\usepackage{array}
\usepackage{xcolor}
\usepackage{listings}
\usepackage{enumitem}
\usepackage{tikz}
\usepackage{pgfplots}
\pgfplotsset{compat=1.18}
\usetikzlibrary{arrows.meta,patterns,calc,decorations.pathmorphing}
\usepackage{cite}
\usepackage[hidelinks]{hyperref}

% --- Noms francais des flottants ---
\renewcommand{\tablename}{Tableau}
\renewcommand{\figurename}{Fig.}
\renewcommand{\lstlistingname}{Listing}
\addto\captionsfrench{\renewcommand{\refname}{Références}}

% --- En-tete de revue commun a tous les TP ---
\newcommand{\improvejournal}{IEEE Transactions on Computational Mechanics and Machine Learning, Vol.~1, No.~1, Juin~2026}

% --- Style code Python (mono, sobre facon IEEE) ---
\definecolor{kw}{HTML}{0033B3}
\definecolor{cm}{HTML}{7A7A7A}
\definecolor{st}{HTML}{067D17}
\lstdefinestyle{py}{%
  language=Python, basicstyle=\ttfamily\footnotesize,
  keywordstyle=\color{kw}, commentstyle=\color{cm}\itshape, stringstyle=\color{st},
  numbers=none, showstringspaces=false, breaklines=true, columns=fullflexible,
  aboveskip=3pt, belowskip=1pt, xleftmargin=2pt,
  literate={é}{{\'e}}1 {è}{{\`e}}1 {à}{{\`a}}1 {ê}{{\^e}}1 {ô}{{\^o}}1 {î}{{\^i}}1 {ç}{{\c c}}1 {É}{{\'E}}1 {_}{{\_}}1}
\lstset{style=py}
